 
%\chapter {Introduction and Motivation}
\section{Introduction} 

Logistics, planning and controlling the efficiency of the flow
and storage activities in the system, is one of the most important 
functions of a company. The design of a distribution system begins 
with the questions of where to locate the facilities and how to 
allocate customers to the selected facilities. 
Furthermore, determining how to distribute goods to customers from
these facilities and how to use vehicles and/or 
distribution resources efficiently are important
decisions that arise in the design of logistics systems. 

In this research, we focus on a class of problems that integrates
the decisions of facility location, vehicle routing, and route assignment 
and seeks to minimize the total cost. We refer to such problems as location 
routing and scheduling problems (LRSPs). The LRSP generalizes and subsumes 
several well-studied problem classes, such as the location and routing problems 
(LRP), the vehicle routing problem (VRP), the multi-depot VRP (MDVRP), and the 
VRP with multiple trips (VRPMT). 

Given a set of candidate facility locations 
and a set of customer locations, the objective of the LRSP is to select a subset 
of the facilities, construct a set of delivery routes, and assign routes to vehicles 
in such a way as to minimize total cost subject to satisfaction of the system's constraints. 
The total cost may include both fixed costs and operating costs of 
facilities and vehicles. 

The LRSP arises in the context of many delivery problems, 
such as gas, oil, food, beverage, bill and newspaper delivery.
In practical applications of the LRS problem, various constraints need to be considered. 
For example, facilities and vehicles may have capacities, the length of routes may
be restricted with time limits, service of customers may include both deliveries 
and pick ups, service to customers might be restricted to occur within time windows, 
different types of vehicles might be  used, customers might have priorities, and some 
data such as demand of customers or travel times might be uncertain. 
In this research, we consider an LRSP with capacity constraints on the 
facilities and on the vehicles, as well as time constraints on the vehicles. 
We assume that all data is deterministic, and the problem is either a pick-up or 
delivery problem. 


\section {Motivation}
The motivation behind the integration of three decisions, facility location, 
vehicle routing and vehicle scheduling is to improve the efficiency of the 
solution and yield more accurate and cost-effective solutions. 
Early in the literature (\citet{Webb} and \citet{Salhi}),
it has been realized that the location decision and the routing 
decision are interdependent when customers demand less-than-truckload quantities.
%and can be served on multiple-stop routes,  
Therefore, if it is possible to serve customers through routes making multiple stops, 
the assumption of single-customer routes in location models will not accurately 
capture the transportation cost. Gas and oil delivery systems, distribution of goods to 
retailers, and newspaper and bill delivery systems are examples of such systems
that require routes with multiple stops.

Location and routing models seek to minimize total cost by simultaneously selecting a 
set of facilities and constructing a set of delivery routes that satisfy the specified 
system constraints. LRPs, however, implicitly assume that each
vehicle covers exactly one route. This may potentially overestimate the number of 
vehicles required and the associated distribution cost. In many cases, it is 
possible to serve multiple routes with a single vehicle, in which case the 
decision of assigning routes to vehicles becomes interdependent with
the location and routing decisions. \citet{Eilon} reports that fixed costs 
are equal to the eighty percent of total vehicle costs. Therefore, models 
integrating the scheduling decision and considering the fixed costs of the 
vehicles may yield more cost efficient solutions. Furthermore, some systems 
require different types of vehicles, which may have different costs 
than the multi-purpose vehicles. In such cases, the number of 
customized vehicles is an important decision for the 
company. Similarly, if the service to the customer 
requires a trained worker who is more expensive than a regular driver, then
the company will want to consider training costs, which are also fixed costs.

Why capturing the costs more accurately is important? 
Major components of logistics cost are warehousing costs including inventory carrying costs 
and transportation costs. Based on estimates for the U.S. in 2002, transportation costs 
were \$577 billion, and inventory carrying costs and warehousing costs were \$298 billion. 
%and logistics administration costs were \$35 billion. 
The total logistics costs were \$910 billion, which was equivalent to 8.7\% of the U.S. gross 
domestic product in 2002 (\citet{MacroSys}). Undeniably, logistics
costs are very important for the U.S. economy as well as for a company. 
%That is why logistics has received such attention by researchers and managers.   
%With an increase in the efficiency of logistics activities, 
%the related costs will decrease and customer satisfaction will increase. 
Since the contribution of logistics activities to the GDP is
large, improvements in the efficiency of the logistics activities will contribute to the 
growth of the U.S. GDP.  Even a small amount of improvement in distribution systems
can result in a substantial amount of savings. 
%Therefore, logistics 
%decisions are very important and require professional decisions and
%efficient solution techniques to optimize the best interests.

Another motivation to consider scheduling decisions together with the location
and routing decisions is to be able include more system constraints to the model. 
In distribution systems in which  the company has a constant fleet size, the number 
of vehicles should be considered in the model. To be able to cover all customers,
it might be necessary to use a vehicle for more than one route. Therefore, the 
decision of scheduling of vehicles must be considered with the decisions of 
location and routing.  
In some systems, there are time constraints for the customers or
the vehicles. Delivery of some items such as perishable foods, newspapers and
bills are time sensitive. In addition, the overtime for the workers might
be expensive, hence the company is not willing to pay overtime driver hours.
Therefore, while designing the model, time limit for the vehicle routes must be
considered.   









\begin{comment}
Designing a distribution network requires decisions regarding 
location, routing and scheduling. Companies must decide
where to locate their facilities, how to serve their customers
on routes and how to use their vehicles and/or 
distribution resources efficiently to optimize the flow on their
distribution networks and to optimize storage costs in their warehouses.  
\end{comment}
\begin{comment}
In this proposal, we study integrated location, routing and scheduling (LRS) problems.
In general terms, the LRS problem integrates the decisions of determining 
the optimal number and locations of facilities, determining an optimal set of
vehicle routes from facilities to customers and determining an optimal assignment
of routes to vehicles  subject to scheduling constraints. The objective is to minimize 
the total system cost, which may include both fixed costs and operating costs of 
facilities and vehicles. LRS problems arise in the context of many delivery problems, 
such as gas and oil delivery, food and beverage delivery and bill and newspaper delivery.
The LRS problem generalizes a number of logistics problems including the location
routing problem (LRP), the multi-depot vehicle routing problem (MDVRP) and the vehicle
routing problem with multiple trips (VRPMT). 

In practical applications of the LRS problem, various constraints need to be considered. 
For example, facilities and vehicles may have capacities, the length of routes may
be restricted by time limits, service of customers may be both deliveries and pick ups, 
service to customers might be restricted to occur within given time windows, 
vehicles might be  different types, customers might have priorities, and some 
data such as demand of customers or travel times might be uncertain. In this 
proposal, we describe our progress to date on developing models and solution
algorithms for a deterministic system with capacitated facilities and vehicles and a 
time restriction for the working hours of the vehicles as scheduling constraint.  
\end{comment}

\begin{comment}
\section{Motivation}
\label{motiv}
Since logistics systems are very complicated and require many
different constraints to be considered, it is almost impossible
to model the exact systems and to optimize them with the current solution 
methods. Therefore, simpler models are developed to approximate the
real systems. The  design of a basic logistics system begins with 
answering questions of (i) where to locate the facilities and 
(ii) how to allocate customers to these facilities. These can be
determined using location-allocation models, which are based on the assumption that 
customers are served individually on out-and-back routes. Location-allocation
models are not 
appropriate for situations in which customers have demands that are 
``less-than-truckload'' and therefore multiple customers can be served on a route 
by one vehicle. Gas and oil delivery systems, distribution of goods to 
retailers, and newspaper and bill delivery systems are examples of such systems
that require routes with multiple stops. 

After determining the location of facilities and the allocation of customers
to the facilities, we can solve a vehicle routing problem for each facility
to determine the customer routes with multiple stops. However,  
the decisions of location and routing are interdependent. 
To represent the distribution system realistically and to find more cost 
efficient solutions, the routing decisions should be included in the location-allocation 
model. Combined location and routing problems consider the multiple stop nature of 
delivery routes in deciding together the facility locations and the optimal vehicle routes. 
 
Location routing models assume that each vehicle covers one route. Ignoring the 
possibility that a vehicle may cover more than one route results in a solution that 
potentially overestimates the number of vehicles required and the associated distribution cost.
\citet{Eilon} reports that fixed costs are equal to the eighty percent of total vehicle costs.
Therefore, a more realistic model  must consider the fixed costs of the 
vehicles to get more cost efficient solutions. Furthermore, some systems require 
different type of vehicles, which may have different costs than the multi-purpose vehicles. 
In such cases, the company also must decide the optimal number of customized vehicles 
necessary to serve its customers. Similarly, if the service to the customer 
requires a trained worker who is more expensive than a regular driver, then
the company will want to consider training costs,  which are fixed costs.
% and want to minimize the total number of trained drivers.    

In distribution systems in which  the company has a constant fleet size, the number 
of vehicles should be considered in the model. To be able to cover all customers,
it might be necessary to use a vehicle for more than one route. Therefore, the 
decision of scheduling of vehicles must be considered with the decisions of 
location and routing.  

Another motivation to consider scheduling decisions together with the location
and routing decisions is the time constraints for the customers or
the vehicles. Delivery of some items such as perishable foods, newspapers and
bills are time sensitive. In addition, the overtime for the workers might
be expensive, hence the company is not willing to pay overtime driver hours.
Therefore, while designing the model, time limit for the vehicle routes must be
considered.   
\end{comment}



%-design of logistic system-->location, routing and scheduling
%-facility location problems
%-fixed facilities vehicle routing
%-location allocation problems
%-imporatnce of scheduling
%examples fleet size is constant
%-save fixed costs 
%-overtime is expensive
%-trained drivers


%A more realistic distribution system must consider the fixed costs of the vehicles
%to get more cost efficient solutions. \citet{Eilon} reports that fixed costs are
%equal to the \%80 of total vehicle costs. Therefore, ignoring the fixed cost of  
%vehicles results in unrealistic cost estimates for the real system. 

%Combined location and routing problems consider 
%the multiple stop nature of delivery routes in deciding the facility locations 
%and the optimal vehicle routes together, but they assume that each vehicle covers one
%route. Ignoring the possibility that a vehicle may cover more than one route results
%in a solution that overestimates the number of vehicles required and the 
%associated distribution cost. The integrated location, routing and scheduling 
%(LRS) problem simultaneously optimizes the decisions of facility location, vehicle
%routing and scheduling.


%[for motivation--> from Petch and Salhi sayfa 70 e bak]

%motivation-most of the real life applications fleet size is constantyou need 
%to assign more than one task to a vehicle.  

%also look at multiobjective LRS paper 

%Since the transportation costs are huge, even a small percentage of 
%reduction in cost may cause large savings. Although in literature, 
%some of the constraints of real-life problems are relaxed to get 
%simpler models, it has been realized that consideration of integrated  
%decisions and combined optimization of these decisions may result more efficient 
%solutions. In the real world, the decisions of locating facilities, 
%determining routes from facilities to customers and scheduling the vehicles
%are interdependent.
